% =====================================================================
%  1bat-ud5-9.tex  —  SESSIÓ 9: Comparar ofertes financeres amb la TAE
%  (sense preàmbul: s'inclou des de main.tex)
% =====================================================================
\horatitol{Sessió 9 — Comparar ofertes financeres amb la TAE}%
{Davant de dues o tres ofertes de crèdit, decidir quina és realment més cara: el TIN i la quota enganyen; només la TAE permet una comparació justa.}

\subsection{Idea clau}

A la sessió 8 vam veure què és la TAE i com es calcula. Avui treballem la competència
\textbf{més transferible} de tot el bloc financer i que aplicaràs acompanyant famílies
endeutades: \textbf{comparar críticament} diverses ofertes de crèdit. La regla és
senzilla però sovint ignorada: \emph{no} es comparen ofertes pel TIN ni per la quota
mensual, sinó per la \textbf{TAE} i pel \textbf{total a retornar}.

\begin{idea}
\textbf{Com es comparen ofertes de crèdit}
\begin{enumerate}[leftmargin=1.6em, itemsep=2pt]
  \item \textbf{Mira la TAE}, no el TIN: la TAE inclou comissions i capitalització; el
        TIN no.
  \item \textbf{Mira el total a retornar} (capital $+$ tots els interessos i
        comissions), no només la quota mensual.
  \item \textbf{Compte amb la lletra petita:} terminis curts, comissions d'obertura,
        assegurances obligatòries.
\end{enumerate}
\textbf{Parany habitual:} una quota mensual baixa pot amagar molts mesos de pagament i
un total caríssim. \textbf{Una TAE més alta vol dir un crèdit més car}, a igualtat de
la resta.
\end{idea}

\subsection{Exemples}

\begin{exemple}[dues ofertes per $\num{1000}\eur$]
Una família necessita $\num{1000}\eur$ i té dues ofertes:
\begin{center}
\renewcommand{\arraystretch}{1.3}
\begin{tabular}{l c c}
\toprule
 & \textbf{Oferta A (microcrèdit)} & \textbf{Oferta B (banc)} \\
\midrule
TIN anunciat & $4\%$ mensual & $9\%$ anual \\
TAE & $60,1\%$ & $9,4\%$ \\
\bottomrule
\end{tabular}
\end{center}
L'oferta A presenta un «$4\%$» que sembla petit, però és \textbf{mensual}: la seva TAE
és $(1,04)^{12}-1\approx 0,601=60,1\%$. L'oferta B té una TAE del $9,4\%$. \textbf{La
B és molt més barata}, encara que el seu número «$9\%$» semblés més gran que el «$4\%$»
de l'A. Comparar pel número de cara enganya; comparar per la TAE no.
\end{exemple}

\begin{exemple}[la quota baixa que surt cara]
Una oferta promociona «només $\num{50}\eur$ al mes!». Sona assumible, però si el
préstec dura $36$ mesos, el total pagat és $36\per\num{50}=\num{1800}\eur$. Si el
capital era de $\num{1200}\eur$, s'han pagat $\num{600}\eur$ d'interessos i comissions.
La quota baixa amagava un \textbf{cost total} alt: per això cal mirar el total a
retornar, no la quota.
\end{exemple}

\subsection{Exercicis resolts}

\begin{exresolt}[quina oferta és més cara]
Dues ofertes per $\num{500}\eur$: l'oferta P té un TIN del $3\%$ mensual; l'oferta Q
té una TAE del $30\%$ anual. Quina és més cara?
\solucio
Calculem la TAE de l'oferta P (que ve donada en mensual):
\[
\text{TAE}_P=(1,03)^{12}-1\approx 1,4258-1=0,4258=42,6\%.
\]
L'oferta Q té TAE $=30\%$. Com $42,6\%>30\%$, l'\textbf{oferta P és més cara}, tot i
que el seu «$3\%$» semblava un número petit.
\resposta{L'oferta P (TAE $42,6\%$) és més cara que la Q (TAE $30\%$).}
\end{exresolt}

\begin{exresolt}[comparar pel total a retornar]
Oferta X: $\num{1000}\eur$ a tornar amb una quota de $\num{95}\eur$ durant $12$ mesos.
Oferta Y: $\num{1000}\eur$ a tornar amb una quota de $\num{60}\eur$ durant $24$ mesos.
Quina fa pagar més en total?
\solucio
\textbf{Oferta X:} $12\per\num{95}=\num{1140}\eur$ en total ($\num{140}\eur$
d'interessos).

\textbf{Oferta Y:} $24\per\num{60}=\num{1440}\eur$ en total ($\num{440}\eur$
d'interessos).

Tot i tenir una quota més baixa ($\num{60}\eur$ contra $\num{95}\eur$), l'\textbf{oferta
Y fa pagar molt més} ($\num{1440}\eur$ contra $\num{1140}\eur$) perquè dura el doble de
temps. La quota baixa enganya.
\resposta{L'oferta Y ($\num{1440}\eur$) surt més cara que la X ($\num{1140}\eur$), tot
i la quota més baixa.}
\end{exresolt}

\subsection{Exercicis per fer}

\begin{exercicis}
\begin{exlist}
  \item Dues ofertes per $\num{800}\eur$: l'oferta A té un TIN del $2,5\%$ mensual;
        l'oferta B, una TAE del $28\%$. Calcula la TAE de l'A i digues quina és més
        cara.
  \item Calcula el \textbf{total a retornar} de cada oferta i digues quina és més cara:
    \begin{enumerate}[label=\alph*), leftmargin=1.6em, topsep=2pt]
      \item quota de $\num{120}\eur$ durant $10$ mesos;
      \item quota de $\num{80}\eur$ durant $18$ mesos.
    \end{enumerate}
  \item Una oferta destaca «TIN $1,5\%$ mensual». Calcula'n la TAE. Et sembla un crèdit
        car o barat comparat amb un préstec personal de banc (TAE típica al voltant del
        $8$–$10\%$)?
  \item \textbf{(Decisió justificada)} Una família et porta dues ofertes: A (TAE
        $45\%$, quota $\num{40}\eur$/mes) i B (TAE $12\%$, quota $\num{55}\eur$/mes).
        Quina els recomanaries i per què? La quota més baixa és sempre la millor opció?
  \item \textbf{(Interpretació)} Per què comparar dos crèdits només per la quota mensual
        pot portar a triar el més car? Quins dos factors amaga la quota?
  \item \textbf{(La meva llista de comprovacions)} Escriu la teva llista personal de
        comprovacions davant d'una oferta de crèdit (almenys tres coses que miraràs
        sempre). Pensa-la com una eina de defensa financera per a tu i per a les
        famílies que assessoraràs.
\end{exlist}
\end{exercicis}

% ---------------------------------------------------------------------
\begin{idea}
\textbf{La meva llista de comprovacions davant d'un crèdit} (per tenir sempre a mà):
\begin{enumerate}[leftmargin=1.6em, itemsep=1pt]
  \item Mirar la \textbf{TAE} (no el TIN ni la quota).
  \item Calcular el \textbf{total a retornar} (quota $\times$ nombre de mesos).
  \item Llegir la \textbf{lletra petita}: comissions, assegurances, terminis.
  \item \textbf{Comparar} sempre amb almenys una altra oferta abans de decidir.
\end{enumerate}
\end{idea}

% ---------------------------------------------------------------------
\fullsolucions{Sessió 9}
\begin{solucions}
\begin{sollist}
  \item TAE de l'oferta A: $(1,025)^{12}-1\approx 1,3449-1=0,3449=34,5\%$. Com
        $34,5\%>28\%$, l'\textbf{oferta A és més cara} que la B.
  \item (a) $10\per\num{120}=\num{1200}\eur$. \quad (b) $18\per\num{80}=\num{1440}\eur$.
        L'oferta (b) és més cara ($\num{1440}\eur$ contra $\num{1200}\eur$), tot i la
        quota més baixa.
  \item TAE $=(1,015)^{12}-1\approx 1,1956-1=0,1956=19,6\%$. És \textbf{car} comparat
        amb un préstec personal de banc (TAE $8$–$10\%$): gairebé el doble, tot i que el
        «$1,5\%$ mensual» semblés mínim.
  \item Resposta oberta. Idea: recomanaria l'\textbf{oferta B} (TAE $12\%$), molt més
        barata que l'A (TAE $45\%$), encara que la quota sigui una mica més alta. La
        quota baixa de l'A no compensa: amb una TAE tan alta, el total a retornar serà
        molt superior. \textbf{No}, la quota més baixa no és sempre la millor opció.
  \item Comparar només per la quota amaga \textbf{el nombre de mesos} (un crèdit amb
        quota baixa però molts mesos pot costar més en total) i \textbf{el cost real}
        (la TAE, amb comissions incloses). Dues quotes baixes poden correspondre a
        crèdits molt diferents.
  \item Resposta oberta. Exemples vàlids: mirar la TAE; calcular el total a retornar;
        llegir la lletra petita (comissions, assegurances); comparar amb una altra
        oferta; desconfiar de la publicitat que destaca el TIN mensual o una quota molt
        baixa.
\end{sollist}
\end{solucions}
